\section{Background and Related Work}
\label{sec:background}
Effective WAF defense against SQLi requires understanding the adversarial landscape well enough to identify and close policy gaps. We review three strands of literature that motivate \emph{family-resolved} coverage reporting rather than relying on headline bypass percentages alone: mutation-based attack research (which reveals where defenses currently fail), parsing discrepancy research (which identifies structural vulnerabilities in WAF architectures), and defensive evaluation frameworks (which articulate what rigorous coverage benchmarks require).

\begin{figure}[t]
  \centering
  \fbox{\parbox{\columnwidth}{\centering Layers: lexical SQLi mutations; URL/query encodings; edge WAF normalization; application parser and ORM; observable HTTP/app signals}}
  \caption{Conceptual layering for SQLi WAF coverage assessment (simplified). Each layer can introduce transformations that affect whether a static rule or model ``sees'' attack tokens---and therefore which layer a defensive rule must address.}
  \label{fig:related-taxonomy}
\end{figure}

\subsection{Mutation-Based SQLi Research and What It Reveals for Defenders}
Adversarial methods search for inputs that preserve attacker intent while evading a detector's decision boundary. SSQLi reports bypass rates up to 97.39\% in its experimental setting using SAC \cite{guan2023ssqli}. GSQLi uses GAN-based mutation \cite{khan2024gsqli}. WAF-A-MoLE applies semantics-preserving fuzzing to stress ML WAFs \cite{demetrio2020wafamole}. From a \emph{defensive} standpoint, these lines of work share a critical structural insight: the attack space is combinatorial---comment insertion, case folding, encoding nests, whitespace and delimiter tricks, and dialect-specific spellings can be composed in ways that interact non-monotonically with regex-based and tokenized defenses: an encoding that defeats one rule may trigger another, or may change how a downstream parser tokenizes a string \cite{mutation-sqli}.

For defenders, the relevant output is therefore not only ``bypass rate'' but \emph{which specific compositions survive}, because survivors identify actionable rule gaps. A mutation family that is uniformly blocked indicates a well-covered area; a family that consistently bypasses identifies a remediation target. Decomposing rates by family converts aggregate statistics into a prioritized repair agenda---which is precisely what this work provides.

\subsection{Parsing Discrepancies and Defensive Scope}
When components disagree on decoding order, charset, or structural boundaries, attackers can sometimes satisfy benign predicates at the edge while delivering malicious semantics to the application \cite{akhavani2025waffled}. WAFFLED's large-scale study underscores that HTTP is not a single abstract datatype but a family of implementations. Our experiments deliberately stabilize parameter binding via remapped seeds so that discrepancy across \emph{different parameter names} is minimized; remaining variance is dominated by representation of the injected value itself. This scoping choice trades ecological completeness (full HTTP discrepancy surface) for interpretability of per-family SQLi mutation outcomes against one policy snapshot---a deliberate methodological decision to enable focused, actionable policy diagnosis.

\subsection{Defensive Frameworks and Coverage Evaluation Requirements}
Hybrid defenses aim to raise evasion cost through better training data \cite{elebiary2026wamm}, shadow models and automated patch generation \cite{wu2025wafbooster}, or input preprocessing \cite{zhang2026bwafsql}. HTTP-Normalizer exemplifies RFC-oriented validation as a discrepancy mitigation \cite{httpnormalizer}. These proposals represent the state of the art in defensive improvement, but they share a common requirement: \emph{paired offensive benchmarks} that tie mutation attempts to family labels and outcome logs. Without such benchmarks, it is impossible to determine whether a reported improvement actually closes high-risk coverage gaps or merely shifts which variants are generated. Our work supplies an open, log-grounded coverage baseline specifically designed to enable evaluation and comparison of defensive rule improvements.

\subsection{Coverage Metrics and Threat of Confounding}
Black-box WAF evaluation is inherently tied to \emph{observable} outcomes: status codes, response bodies, and timing. Different probes may map the same underlying vulnerability to different HTTP observables (e.g., generic 500 versus structured 422). Any operational definition of ``policy gap'' therefore embeds epistemic commitments; transparent scripting and published logs mitigate but do not eliminate this issue. Section~\ref{sec:methodology} states our audit probe's outcome labels explicitly; Section~\ref{sec:discussion} revisits threats to validity and their implications for interpreting and acting on the coverage findings.
